% Options for packages loaded elsewhere
\PassOptionsToPackage{unicode}{hyperref}
\PassOptionsToPackage{hyphens}{url}
\documentclass[
  a4paper,11pt]{article}
\usepackage{xcolor}
\usepackage[margin=1in]{geometry}
\usepackage{amsmath,amssymb}
\usepackage{graphicx}
\setcounter{secnumdepth}{-\maxdimen} % remove section numbering
\usepackage{iftex}
\ifPDFTeX
  \usepackage[T1]{fontenc}
  \usepackage[utf8]{inputenc}
  \usepackage{textcomp}
\else
  \usepackage{unicode-math}
  \defaultfontfeatures{Scale=MatchLowercase}
  \defaultfontfeatures[\rmfamily]{Ligatures=TeX,Scale=1}
\fi
\usepackage{lmodern}
\ifPDFTeX\else
  \setmainfont[]{Times New Roman}
\fi
\IfFileExists{upquote.sty}{\usepackage{upquote}}{}
\IfFileExists{microtype.sty}{%
  \usepackage[]{microtype}
  \UseMicrotypeSet[protrusion]{basicmath}
}{}
\makeatletter
\@ifundefined{KOMAClassName}{%
  \IfFileExists{parskip.sty}{%
    \usepackage{parskip}
  }{%
    \setlength{\parindent}{0pt}
    \setlength{\parskip}{6pt plus 2pt minus 1pt}}
}{% if KOMA class
  \KOMAoptions{parskip=half}}
\makeatother
\usepackage{longtable,booktabs,array}
\newcounter{none}
\usepackage{calc}
\usepackage{etoolbox}
\makeatletter
\patchcmd\longtable{\par}{\if@noskipsec\mbox{}\fi\par}{}{}
\makeatother
\IfFileExists{footnotehyper.sty}{\usepackage{footnotehyper}}{\usepackage{footnote}}
\makesavenoteenv{longtable}
\usepackage{bookmark}
\IfFileExists{xurl.sty}{\usepackage{xurl}}{}
\urlstyle{same}
\usepackage{hyperref}
\hypersetup{
  hidelinks,
  pdfcreator={LaTeX via pandoc}}
\setlength{\emergencystretch}{3em}
\providecommand{\tightlist}{%
  \setlength{\itemsep}{0pt}\setlength{\parskip}{0pt}}
\usepackage{amsthm}
\newtheorem{theorem}{Theorem}[section]
\newtheorem{proposition}[theorem]{Proposition}
\newtheorem{lemma}[theorem]{Lemma}
\newtheorem{corollary}[theorem]{Corollary}
\theoremstyle{definition}
\newtheorem{definition}[theorem]{Definition}
\theoremstyle{remark}
\newtheorem{remark}[theorem]{Remark}
\newtheorem{observation}[theorem]{Observation}
\begin{document}
\section{Set-Theoretic Structure and Combinatorial Properties of the
6-Dimensional
Hypercube}\label{set-theoretic-structure-and-combinatorial-properties-of-the-6-dimensional-hypercube}

\subsection{------ Spin Classification, Set Enumeration, and Sign
Product
------}\label{spin-classification-set-enumeration-and-sign-product}

\textbf{Author}: Noriaki Kihara (WF System Co., Ltd.) \textbf{Date}:
April 2026 \textbf{DOI}: 10.5281/zenodo.19646652

\begin{center}\rule{0.5\linewidth}{0.5pt}\end{center}

\subsection{0. Purpose of This Paper}\label{purpose-of-this-paper}

This paper defines a 6-dimensional structure obtained by adding a
discrete label axis \(Q\) to a 5-dimensional hypercube, and describes
the combinatorial properties derived from its set-theoretic structure.

In the main body (§1--§8), we assert only the mathematical structure of
this model and the consequences derived purely combinatorially
therefrom. The correspondence with the Standard Model is presented in §9
as ``an example of interpretation,'' but this is not a claim of this
paper. It is merely one example showing how the set-theoretic structure
of the model may be mapped to the particle classification of the
Standard Model.

\textbf{On the external assumptions of this model}: This paper takes the
5-dimensional hypercube as its starting point. Why 5 dimensions, and why
certain axes among the 5 are later treated as ``spatial axes,'' are
\textbf{external assumptions of this model} and are not justified within
the model itself. This dimensional setting is adopted as a
correspondence to the empirical fact that the observed spacetime
consists of 3 spatial dimensions + 1 time dimension.

\begin{center}\rule{0.5\linewidth}{0.5pt}\end{center}

\subsection{1. Definition of the 6-Dimensional
Structure}\label{definition-of-the-6-dimensional-structure}

\subsubsection{1.1 Definition}\label{definition}

The 6-dimensional structure treated in this paper is defined as follows.

\begin{itemize}
\tightlist
\item
  \textbf{Geometric structure}: The combinatorial structure of the
  5-dimensional hypercube (vertices, edges, faces, 3-cells, 4-facets,
  5-cell body)
\item
  \textbf{Label axis}: An 8-valued discrete label axis
  \(Q \in \{0, 1, 2, 3, 4, 5, 6, 7\}\)
\end{itemize}

The vertex coordinates of the 5-dimensional hypercube are denoted by
\((x_1, x_2, x_3, x_4, x_5)\), where each coordinate \textbf{takes
arbitrary real values}. When a specific proposition imposes conditions
on coordinate values within the set-theoretic structure, such conditions
are declared locally within that proposition.

The 6th axis \(Q\) is not a geometric coordinate but a discrete label
value associated with each state. Its range is the 8 values
\(\{0, 1, 2, 3, 4, 5, 6, 7\}\). These 8 values possess a 3-bit
structure, encoded as \(Q = 4c_1 + 2c_2 + c_3\) (\(c_i \in \{0, 1\}\)).
The addition of the \(Q\) axis does not change the combinatorial
structure (number of vertices, edges, faces, etc.) of the hypercube.

\subsubsection{\texorpdfstring{1.2 Element Counts of the
\(k\)-Dimensional
Hypercube}{1.2 Element Counts of the k-Dimensional Hypercube}}\label{element-counts-of-the-k-dimensional-hypercube}

Let \(f_j(k)\) denote the number of \(j\)-dimensional faces
(\(j\)-faces) of the \(k\)-dimensional hypercube. We present the element
counts starting from low dimensions so that the reader can verify the
general formula.

\textbf{Number of vertices \(f_0(k)\)}: As the combinatorial count
distinguishing the 2 states at both ends of each axis:

{\def\LTcaptype{none} % do not increment counter
\begin{longtable}[]{@{}ccccccc@{}}
\toprule\noalign{}
\(k\) & 0 & 1 & 2 & 3 & 4 & 5 \\
\midrule\noalign{}
\endhead
\bottomrule\noalign{}
\endlastfoot
\(f_0(k)\) & 1 & 2 & 4 & 8 & 16 & 32 \\
\end{longtable}
}

\[f_0(k) = 2^k\]

\textbf{Number of edges \(f_1(k)\)}: Each edge extends along 1 axis
direction and is determined by the position at both ends of the
remaining \((k-1)\) axes:

{\def\LTcaptype{none} % do not increment counter
\begin{longtable}[]{@{}ccccccc@{}}
\toprule\noalign{}
\(k\) & 0 & 1 & 2 & 3 & 4 & 5 \\
\midrule\noalign{}
\endhead
\bottomrule\noalign{}
\endlastfoot
\(f_1(k)\) & 0 & 1 & 4 & 12 & 32 & 80 \\
\end{longtable}
}

\[f_1(k) = k \cdot 2^{k-1}\]

Verification: For \(k=3\), \(3 \cdot 4 = 12\) (12 edges of a cube); for
\(k=5\), \(5 \cdot 16 = 80\).

\textbf{Number of square faces \(f_2(k)\)}: Faces spanned by 2 axes:

{\def\LTcaptype{none} % do not increment counter
\begin{longtable}[]{@{}ccccccc@{}}
\toprule\noalign{}
\(k\) & 0 & 1 & 2 & 3 & 4 & 5 \\
\midrule\noalign{}
\endhead
\bottomrule\noalign{}
\endlastfoot
\(f_2(k)\) & 0 & 0 & 1 & 6 & 24 & 80 \\
\end{longtable}
}

\[f_2(k) = \binom{k}{2} \cdot 2^{k-2}\]

Verification: For \(k=3\), \(3 \cdot 2 = 6\) (6 faces of a cube); for
\(k=5\), \(10 \cdot 8 = 80\).

\textbf{Number of cubic cells (3-cells) \(f_3(k)\)}: Cubic cells spanned
by 3 axes:

{\def\LTcaptype{none} % do not increment counter
\begin{longtable}[]{@{}ccccccc@{}}
\toprule\noalign{}
\(k\) & 0 & 1 & 2 & 3 & 4 & 5 \\
\midrule\noalign{}
\endhead
\bottomrule\noalign{}
\endlastfoot
\(f_3(k)\) & 0 & 0 & 0 & 1 & 8 & 40 \\
\end{longtable}
}

\[f_3(k) = \binom{k}{3} \cdot 2^{k-3}\]

Verification: For \(k=4\), \(4 \cdot 2 = 8\); for \(k=5\),
\(10 \cdot 4 = 40\).

\textbf{Number of 4-facets \(f_4(k)\)}: 4-dimensional cells spanned by 4
axes:

{\def\LTcaptype{none} % do not increment counter
\begin{longtable}[]{@{}ccccccc@{}}
\toprule\noalign{}
\(k\) & 0 & 1 & 2 & 3 & 4 & 5 \\
\midrule\noalign{}
\endhead
\bottomrule\noalign{}
\endlastfoot
\(f_4(k)\) & 0 & 0 & 0 & 0 & 1 & 10 \\
\end{longtable}
}

\[f_4(k) = \binom{k}{4} \cdot 2^{k-4}\]

Verification: For \(k=5\), \(5 \cdot 2 = 10\).

\textbf{Number of 5-cell bodies \(f_5(k)\)}: For \(k=5\):

\[f_5(5) = \binom{5}{5} \cdot 2^0 = 1\]

\subsubsection{General Formula}\label{general-formula}

Summarizing the above, the number of \(j\)-dimensional faces of the
\(k\)-dimensional hypercube is given by

\[f_j(k) = \binom{k}{j} \cdot 2^{k-j} \qquad (0 \leq j \leq k)\]

This corresponds to choosing \(j\) axes from \(k\) to span a \(j\)-face
(\(\binom{k}{j}\) ways), and specifying which end of each of the
remaining \((k-j)\) axes the face is located at (\(2^{k-j}\) ways).

\textbf{Verification (Euler characteristic)}:
\(\sum_{j=0}^{k} (-1)^j f_j(k) = (2-1)^k = 1\), which is consistent with
the hypercube being a contractible space.

\subsubsection{\texorpdfstring{Values Used in This Paper
(\(k=5\))}{Values Used in This Paper (k=5)}}\label{values-used-in-this-paper-k5}

{\def\LTcaptype{none} % do not increment counter
\begin{longtable}[]{@{}clc@{}}
\toprule\noalign{}
\(j\) & Element & \(f_j(5)\) \\
\midrule\noalign{}
\endhead
\bottomrule\noalign{}
\endlastfoot
0 & Vertices & 32 \\
1 & Edges & 80 \\
2 & Square faces & 80 \\
3 & Cubic cells (3-cells) & 40 \\
4 & 4-facets & 10 \\
5 & 5-cell body & 1 \\
\end{longtable}
}

\subsubsection{1.3 Axial Symmetry and Definition of Set
Codes}\label{axial-symmetry-and-definition-of-set-codes}

\paragraph{1.3.1 Complete Symmetry of the 5
Axes}\label{complete-symmetry-of-the-5-axes}

\textbf{Proposition 1.1 (Complete symmetry of the 5 axes)}. The 5
geometric axes of the 5-dimensional hypercube are all equivalent in
their combinatorial structure. Any permutation of axes gives an
automorphism of the hypercube, and the 5 axes are indistinguishable with
respect to the element counts \(f_j(k)\) of §1.2 and subsequent
geometric propositions.

\textbf{Proof}. The automorphism group of the combinatorial structure of
the \(k\)-dimensional hypercube is the hyperoctahedral group \(B_k\),
which contains all axis permutations. Therefore the 5 axes are
combinatorially completely symmetric. \(\square\)

\paragraph{1.3.2 Classification of Axes}\label{classification-of-axes}

\textbf{Definition 1.2 (Notational labels for axes)}. The following
notational labels are assigned to the 5 axes
\(x_1, x_2, x_3, x_4, x_5\):

\begin{itemize}
\tightlist
\item
  \(x_1, x_2, x_3\) are denoted \(x, y, z\)
\item
  \(x_4\) is denoted \(t\)
\item
  \(x_5\) is denoted \(R\)
\end{itemize}

By Proposition 1.1, the 5 axes are equivalent in combinatorial
structure, so the assignment of labels to axes is merely a definitional
choice. These labels are used as notational prerequisites for the
definition of set codes (§1.3.3) and the definition of spin (§2).

\textbf{\(x, y, z, t, R\) are purely notational labels and do not imply
any geometric distinction among the axes.} The operation of assigning
physical interpretations to each axis is carried out in §9 (an example
of interpretation), which is an interpretation from outside the model
and is independent of the geometric structure of §1.

\paragraph{1.3.3 Rigorous Definition of Set
Codes}\label{rigorous-definition-of-set-codes}

Each geometric axis component \(x_i\) (\(i = 1, \ldots, 5\)) is a symbol
representing the \textbf{range constraint} that the coordinate value
along that axis must satisfy. This paper uses the following 5 types of
symbols.

{\def\LTcaptype{none} % do not increment counter
\begin{longtable}[]{@{}cll@{}}
\toprule\noalign{}
Symbol & Meaning & Range of coordinate values \\
\midrule\noalign{}
\endhead
\bottomrule\noalign{}
\endlastfoot
\(+\) & Positive & \(x_i > 0\) \\
\(-\) & Negative & \(x_i < 0\) \\
\(0\) & Origin & \(x_i = 0\) \\
\(\pm\) & Nonzero & \(x_i \neq 0\) (positive or negative) \\
\(*\) & Arbitrary & \(x_i\) is any real number \\
\end{longtable}
}

For the \(Q\) axis, symbols representing the range of values are used
similarly.

{\def\LTcaptype{none} % do not increment counter
\begin{longtable}[]{@{}
  >{\centering\arraybackslash}p{(\linewidth - 4\tabcolsep) * \real{0.1235}}
  >{\raggedright\arraybackslash}p{(\linewidth - 4\tabcolsep) * \real{0.2469}}
  >{\raggedright\arraybackslash}p{(\linewidth - 4\tabcolsep) * \real{0.6296}}@{}}
\toprule\noalign{}
\begin{minipage}[b]{\linewidth}\centering
Symbol
\end{minipage} & \begin{minipage}[b]{\linewidth}\raggedright
Meaning
\end{minipage} & \begin{minipage}[b]{\linewidth}\raggedright
Value of \(Q\)
\end{minipage} \\
\midrule\noalign{}
\endhead
\bottomrule\noalign{}
\endlastfoot
\(0\) & Zero label & \(Q = 0\) \\
\(1\)--\(7\) & Specified discrete value & \(Q\) is the specified integer
value \\
\(\bar{Q}\) & Anti-label & \(Q \to (Q + 4) \bmod 8\) (\(Q = 0\) is
invariant) \\
\(*\) & Arbitrary & \(Q \in \{0, 1, 2, 3, 4, 5, 6, 7\}\) (any value) \\
\end{longtable}
}

The 8 values of \(Q\) have a 3-bit structure \((c_1, c_2, c_3)\). The
\(c_1\) bit (the most significant bit of \(Q\): whether \(Q \geq 4\))
and the 2 bits \(c_2 c_3\) (\(Q \bmod 4\)) encode two independent
discrete degrees of freedom. The physical interpretation is treated in
§9.2.

\textbf{Remark}. Set codes are not coordinate values themselves, but
symbols specifying the \textbf{conditions} that coordinate values must
satisfy. As stated in §1.1, coordinate values in this model can take
arbitrary real values, but in specific propositions or set-theoretic
structures, the range constraints listed above are imposed on the
coordinate values of the relevant axes.

\textbf{Example}. The set \((+, +, 0, 0, 0, 0)\) refers to the
collection of states satisfying
\(x_1 > 0, x_2 > 0, x_3 = 0, x_4 = 0, x_5 = 0, Q = 0\).

\begin{center}\rule{0.5\linewidth}{0.5pt}\end{center}

\subsection{2. Definition of Spin}\label{definition-of-spin}

\textbf{Definition 2.1 (Spin)}. For a set \((x, y, z, t, R, Q)\), the
spin \(s\) is defined as follows.

\textbf{Fermion/boson determination}: Let \(k\) denote the number of
axes among the 4 axes \(xyzt\) whose set code is not \(0\).

\begin{itemize}
\tightlist
\item
  When \(k = 2\): \textbf{Fermionic type}
\item
  When \(k \neq 2\): \textbf{Bosonic type}
\end{itemize}

\textbf{Determination of the spin quantum number}: The number of axes
\(n\) contributing to the spin is defined as follows. For the \(Q\)
axis, only \(Q\) label transitions (§3) contribute to the spin; a static
\(Q\) label value (even if \(Q \neq 0\), when no transition is involved)
is not counted toward \(n\).

\begin{itemize}
\tightlist
\item
  Fermionic type: \(n\) is the number of nonzero elements among
  \(\{R,\, Q\text{-transition}\}\). \(s = n + 1/2\)
\item
  Bosonic type: \(n\) is the number of nonzero elements among
  \(\{t,\, R,\, Q\text{-transition}\}\). \(s = n\)
\end{itemize}

The possible values of \(n\) are \(0, 1, 2\) for the fermionic type and
\(0, 1, 2, 3\) for the bosonic type, so the possible values of spin
\(s\) in this model are as follows.

{\def\LTcaptype{none} % do not increment counter
\begin{longtable}[]{@{}ccc@{}}
\toprule\noalign{}
\(n\) & Boson \(s = n\) & Fermion \(s = n + 1/2\) \\
\midrule\noalign{}
\endhead
\bottomrule\noalign{}
\endlastfoot
0 & \(s = 0\) & \(s = 1/2\) \\
1 & \(s = 1\) & \(s = 3/2\) \\
2 & \(s = 2\) & \(s = 5/2\) \\
3 & \(s = 3\) & ------ \\
\end{longtable}
}

\textbf{Remark 1}. The 4 axes \(xyzt\) are completely equivalent in
their initial symmetry, and the process by which one axis separates as
``temporal'' is interpreted as spontaneous symmetry breaking (§9.1). The
fermionic determination condition \(k = 2\) corresponds to the case
where 2 out of 4 equivalent axes are nonzero, and is not restricted to
the 3 axes \(xyz\).

\textbf{Remark 2}. This definition is a combinatorial definition
internal to this model. The terms ``boson'' and ``fermion'' used here
are terminology within this model; the correspondence to physical spin
angular momentum, statistics, \((-1)^{2s}\) transformation properties,
etc. is treated in §9 (an example of interpretation).

\textbf{Remark 3}. The axes counted toward \(n\) differ between the
fermionic and bosonic types (fermions do not count \(t\)). This
asymmetry is based on the following geometric motivation. In the
fermionic type, exactly 2 out of the 4 axes \(xyzt\) are nonzero, and
the choice of these 2 axes already carries \(\binom{4}{2} = 6\) degrees
of freedom. Among these, the combinations including the \(t\) axis (3
combinations) and those not including it (3 combinations) have both
already participated in the \(k = 2\) determination. Therefore the
degree of freedom of \(t\) has been consumed in determining the
fermionic nature, and is reflected in the \(1/2\) of the spin quantum
number \(s = n + 1/2\). As a mathematically equivalent notation, using
the number \(n'\) of nonzero elements among
\(\{t, R, Q\text{-transition}\}\), one can also write
\(s = (n' - 1) + 1/2 = n' - 1/2\) (for fermions with nonzero \(t\),
\(n' = n + 1\)), but this paper adopts the asymmetric form above for
clarity of definition.

\begin{center}\rule{0.5\linewidth}{0.5pt}\end{center}

\subsection{\texorpdfstring{3. Enumeration of \(Q\) Label
Transitions}{3. Enumeration of Q Label Transitions}}\label{enumeration-of-q-label-transitions}

We consider transitions among the 7 nonzero values
\(\{1, 2, 3, 4, 5, 6, 7\}\) of the 8 values of the \(Q\) axis.
Transitions are described by ordered pairs
\((Q_\text{initial}, Q_\text{final})\).

Based on the 3-bit structure \((c_1, c_2, c_3)\), transitions are
classified into 2 types.

\textbf{\(c_2 c_3\) transitions}: Transitions in which the lower 2 bits
\(c_2 c_3\) change. The ordered pairs among the 3 nonzero values
\(\{01, 10, 11\}\) of \(c_2 c_3\) number \(3 \times 3 = 9\) (including 3
diagonal elements).

\textbf{\(c_1\) transitions}: Transitions in which only the most
significant bit \(c_1\) changes. These are transitions between
\(c_1 \in \{0, 1\}\), and for each of the 4 values of \(c_2 c_3\)
(\(00, 01, 10, 11\)), there exist 2 directions: \(0 \to 1\) and
\(1 \to 0\). Therefore the number of \(c_1\) transitions is
\(4 \times 2 = 8\).

\textbf{Remark}. The operation of extracting independent transition
states from the 9 \(c_2 c_3\) transitions (e.g., removing the linear
combination of diagonal elements) depends on representation-theoretic
structure and is therefore not performed in the main body. The physical
interpretation of the number of independent states is treated in §9.3.

\begin{center}\rule{0.5\linewidth}{0.5pt}\end{center}

\subsection{4. Enumeration of Spin 1
Sets}\label{enumeration-of-spin-1-sets}

Spin 1 sets are, by Definition 2.1, bosonic type (\(k \neq 2\) among
\(xyzt\)) with \(n = 1\). We consider the case where all \(xyzt\) axes
are \(0\) (\(k = 0\)).

The states with exactly 1 nonzero element among
\(\{t, R, Q\text{-transition}\}\) are classified as follows.

{\def\LTcaptype{none} % do not increment counter
\begin{longtable}[]{@{}
  >{\raggedright\arraybackslash}p{(\linewidth - 4\tabcolsep) * \real{0.2273}}
  >{\raggedright\arraybackslash}p{(\linewidth - 4\tabcolsep) * \real{0.4545}}
  >{\centering\arraybackslash}p{(\linewidth - 4\tabcolsep) * \real{0.3182}}@{}}
\toprule\noalign{}
\begin{minipage}[b]{\linewidth}\raggedright
Nonzero axis
\end{minipage} & \begin{minipage}[b]{\linewidth}\raggedright
Set form
\end{minipage} & \begin{minipage}[b]{\linewidth}\centering
Number of states
\end{minipage} \\
\midrule\noalign{}
\endhead
\bottomrule\noalign{}
\endlastfoot
\(t\) & \((0, 0, 0, \pm, 0, 0)\) & 2 (2 choices of sign \(\pm\)) \\
\(R\) & \((0, 0, 0, 0, \pm, 0)\) & 2 (2 choices of sign \(\pm\)) \\
\(Q\)-transition (\(c_2 c_3\)) &
\((0, 0, 0, 0, 0, Q_\text{initial} \to Q_\text{final})\) & 9 (by §3) \\
\(Q\)-transition (\(c_1\)) &
\((0, 0, 0, 0, 0, Q_\text{initial} \to Q_\text{final})\) & 8 (by §3) \\
\end{longtable}
}

\textbf{Total}: \(t\)-fixed 2 states + \(R\)-fixed 2 states +
\(c_2 c_3\) transition 9 states + \(c_1\) transition 8 states =
\textbf{21 states}.

When only \(t\) among \(xyzt\) is nonzero (\(k = 1\)), the state is also
bosonic type and includes spin 1 states where \(t\) contributes to the
count of \(n\).

\textbf{Remark}. The extraction of independent states from the 9
\(c_2 c_3\) transition states (reduction \(9 \to 8\)) and the
correspondence to physical gauge bosons are treated in §9.3--§9.4.

\begin{center}\rule{0.5\linewidth}{0.5pt}\end{center}

\subsection{5. Enumeration of Spin 2
Sets}\label{enumeration-of-spin-2-sets}

Spin 2 sets are, by Definition 2.1, bosonic type with \(n = 2\), i.e., 2
of the 3 elements \(\{t, R, Q\text{-transition}\}\) are nonzero, with
the number of nonzero axes among \(xyzt\) satisfying \(k \neq 2\).

For \(k = 0\) (all \(xyzt\) axes are \(0\)), the combinations of 2
nonzero elements are the following 3 cases.

{\def\LTcaptype{none} % do not increment counter
\begin{longtable}[]{@{}ll@{}}
\toprule\noalign{}
Nonzero 2 elements & Set form \\
\midrule\noalign{}
\endhead
\bottomrule\noalign{}
\endlastfoot
\(t, R\) & \((0, 0, 0, \pm, \pm, 0)\) \\
\(t, Q\)-transition & \((0, 0, 0, \pm, 0, Q_i \to Q_j)\) \\
\(R, Q\)-transition & \((0, 0, 0, 0, \pm, Q_i \to Q_j)\) \\
\end{longtable}
}

Therefore, the number of types of spin 2 sets in this model is
\textbf{3}.

\begin{center}\rule{0.5\linewidth}{0.5pt}\end{center}

\subsection{6. Enumeration of Spin 3
Sets}\label{enumeration-of-spin-3-sets}

Spin 3 sets are, by Definition 2.1, bosonic type with \(n = 3\), i.e.,
all 3 elements \(\{t, R, Q\text{-transition}\}\) are nonzero, with the
number of nonzero axes among \(xyzt\) satisfying \(k \neq 2\).

{\def\LTcaptype{none} % do not increment counter
\begin{longtable}[]{@{}ll@{}}
\toprule\noalign{}
Nonzero 3 elements & Set form \\
\midrule\noalign{}
\endhead
\bottomrule\noalign{}
\endlastfoot
\(t, R, Q\)-transition & \((0, 0, 0, \pm, \pm, Q_i \to Q_j)\) \\
\end{longtable}
}

Therefore, the number of types of spin 3 sets in this model is
\textbf{1}.

\begin{center}\rule{0.5\linewidth}{0.5pt}\end{center}

\subsection{7. Sign Product of Nonzero
Axes}\label{sign-product-of-nonzero-axes}

\textbf{Definition 7.1 (Sign product)}. For a bosonic set, when the axes
among \(t, R\) whose set code is not \(0\) each take a sign of \(+\) or
\(-\), the product of these signs is defined as the \textbf{sign
product} \(P\).

\[P = (-1)^{|\{i \in \{t, R\} : \sigma_i < 0\}|}\]

That is, \(P = +1\) if the number of axes among \(t, R\) with negative
sign is even, and \(P = -1\) if odd.

\textbf{Proposition 7.1 (Classification of \(P\) for each spin type)}.
The possible values of \(P\) for each bosonic type are as follows.

{\def\LTcaptype{none} % do not increment counter
\begin{longtable}[]{@{}
  >{\raggedright\arraybackslash}p{(\linewidth - 4\tabcolsep) * \real{0.3385}}
  >{\raggedright\arraybackslash}p{(\linewidth - 4\tabcolsep) * \real{0.3077}}
  >{\raggedright\arraybackslash}p{(\linewidth - 4\tabcolsep) * \real{0.3538}}@{}}
\toprule\noalign{}
\begin{minipage}[b]{\linewidth}\raggedright
Spin type
\end{minipage} & \begin{minipage}[b]{\linewidth}\raggedright
Nonzero \(\{t, R\}\)
\end{minipage} & \begin{minipage}[b]{\linewidth}\raggedright
Possible values of \(P\)
\end{minipage} \\
\midrule\noalign{}
\endhead
\bottomrule\noalign{}
\endlastfoot
\(s = 0\) & None & \(P = +1\) (unique) \\
\(s = 1\), \(t\) type & \(t\) only & \(P = +1\) or \(-1\) \\
\(s = 1\), \(R\) type & \(R\) only & \(P = +1\) or \(-1\) \\
\(s = 1\), \(Q\)-transition type & None & \(P = +1\) (unique) \\
\(s = 2\), \(tR\) type & Both \(t, R\) & \(P = +1\) or \(-1\) \\
\(s = 2\), \(tQ\) type & \(t\) only & \(P = +1\) or \(-1\) \\
\(s = 2\), \(RQ\) type & \(R\) only & \(P = +1\) or \(-1\) \\
\(s = 3\), \(tRQ\) type & Both \(t, R\) & \(P = +1\) or \(-1\) \\
\end{longtable}
}

\textbf{Proof}. When both \(t, R\) are zero, \(P = +1\) (0 negative
axes). When only 1 is nonzero, \(\sigma_i = \pm\) gives \(P = \pm 1\).
When both are nonzero, the 4 configurations \((+)(+) \to P = +1\),
\((+)(-) \to P = -1\), \((-)(+) \to P = -1\), \((-)(-) \to P = +1\) are
possible, giving \(P = \pm 1\). \(\square\)

The physical interpretation of \(P\) (correspondence with the
directionality of forces) is treated in §9 (an example of
interpretation).

\textbf{Remark}. The orientation structure of signed areas \(M(\sigma)\)
for square faces containing the \(R\) axis, and the rigorous definition
of conserved quantities based thereon, exceed the scope of this paper
and will be treated in a separate paper. Interpretive speculation on the
relation between \(M(\sigma)\) and the mass metric is given in §9.11.1.

\begin{center}\rule{0.5\linewidth}{0.5pt}\end{center}

\subsection{8. Conclusion of the Main
Body}\label{conclusion-of-the-main-body}

In this paper, we defined a 6-dimensional structure obtained by adding
an 8-valued discrete label axis \(Q\) to a 5-dimensional hypercube, and
derived the following combinatorial properties from its set-theoretic
structure.

\begin{enumerate}
\def\labelenumi{\arabic{enumi}.}
\tightlist
\item
  \textbf{Element counts of the hypercube} (§1.2): The number of
  \(j\)-dimensional faces of the \(k\)-dimensional hypercube is
  \(f_j(k) = \binom{k}{j} \cdot 2^{k-j}\). For \(k = 5\): 32 vertices,
  80 edges, 80 square faces, 40 cubic cells, 10 hypercubes.
\item
  \textbf{Complete symmetry of the 5 axes} (Proposition 1.1): The 5
  geometric axes are all equivalent in combinatorial structure, and any
  axis permutation gives an automorphism of the hypercube.
\item
  \textbf{Notational labels for axes} (Definition 1.2): The notational
  labels \(x, y, z, t, R\) are assigned to the 5 axes.
\item
  \textbf{Rigorous definition of set codes} (§1.3.3): Range constraints
  on each axis are expressed using the 5 symbols \(+, -, 0, \pm, *\).
  The \(Q\) axis has 8 values (3-bit structure).
\item
  \textbf{Definition of spin} (Definition 2.1): The fermionic type
  (\(k = 2\)) and bosonic type (\(k \neq 2\)) are determined by the
  number \(k\) of nonzero axes among the 4 axes \(xyzt\), and the spin
  is defined as \(s = n + 1/2\) or \(s = n\) from the number \(n\) of
  nonzero elements among \(\{R, Q\text{-transition}\}\) (fermionic type)
  or \(\{t, R, Q\text{-transition}\}\) (bosonic type).
\item
  \textbf{Enumeration of \(Q\) label transitions} (§3): \(c_2 c_3\)
  transitions (9 cases) and \(c_1\) transitions (8 cases) based on the
  3-bit structure.
\item
  \textbf{Enumeration of spin 1 sets} (§4): \(t\)-fixed 2 states +
  \(R\)-fixed 2 states + \(c_2 c_3\) transition 9 states + \(c_1\)
  transition 8 states = 21 states.
\item
  \textbf{3 types of spin 2 sets} (§5): \(tR, tQ, RQ\) --- 3 types.
\item
  \textbf{1 type of spin 3 sets} (§6): \(tRQ\) --- 1 type.
\item
  \textbf{Sign product} (Definition 7.1, Proposition 7.1): The sign
  product \(P = (-1)^{|\{i : \sigma_i < 0\}|}\) is defined from the
  number of negative-sign axes among \(t, R\), and the possible values
  of \(P\) (unique or both values) are classified for each spin type.
\end{enumerate}

All of these are results derived purely from combinatorics and geometry
under the external assumptions of this model (adoption of the
5-dimensional hypercube, addition of the \(Q\) axis).

\begin{center}\rule{0.5\linewidth}{0.5pt}\end{center}

\subsection{9. An Example of Interpretation: Correspondence with the
Standard
Model}\label{an-example-of-interpretation-correspondence-with-the-standard-model}

\textbf{The content of this section is not a claim of this paper but an
example of interpretation. It does not attempt to physically derive the
Standard Model from this model.} This section presents one example of
how the 6-dimensional structure defined and derived in §1--§8 and its
mathematical properties may be mapped to the particle classification of
the Standard Model.

The following correspondences are not justified from within this model
and are \textbf{choices of interpretation from outside}. If an
alternative mapping becomes available in the future, it can be added in
parallel with this section.

\subsubsection{9.1 Assignment of Physical
Labels}\label{assignment-of-physical-labels}

As stated in §1.3.2, the 5 axes are all equivalent in combinatorial
structure; however, for the purpose of correspondence with the Standard
Model, the following labels are assigned.

\begin{itemize}
\tightlist
\item
  \(x_1, x_2, x_3 \to xyz\): 3-dimensional spatial axes
\item
  \(x_4 \to t\): Time axis
\item
  \(x_5 \to R\): Spatial scale axis (an axis connecting the subjective
  coordinate system with the spatial scale of the 6-dimensional
  structure)
\item
  \(Q\): Gauge symmetry axis (discrete encoding of
  \(\mathrm{SU}(3)_C \times \mathrm{SU}(2)_L\))
\end{itemize}

\textbf{Initial symmetry of \(xyzt\)}: By Proposition 1.1, the 4 axes
\(xyzt\) are completely equivalent in the initial state. The process by
which one of these axes separates as ``time'' in observation is
interpreted as spontaneous symmetry breaking. This symmetry breaking
gives rise to the distinction between charged particles (\(t\) axis
nonzero) and neutral particles (\(t\) axis zero).

\subsubsection{\texorpdfstring{9.2 Correspondence Between the 8 Values
of the \(Q\) Axis and Gauge
Symmetry}{9.2 Correspondence Between the 8 Values of the Q Axis and Gauge Symmetry}}\label{correspondence-between-the-8-values-of-the-q-axis-and-gauge-symmetry}

The 8 values of the \(Q\) axis consist of 3 bits \((c_1, c_2, c_3)\).
These 3 bits encode 2 gauge degrees of freedom.

{\def\LTcaptype{none} % do not increment counter
\begin{longtable}[]{@{}
  >{\centering\arraybackslash}p{(\linewidth - 8\tabcolsep) * \real{0.0375}}
  >{\centering\arraybackslash}p{(\linewidth - 8\tabcolsep) * \real{0.2250}}
  >{\centering\arraybackslash}p{(\linewidth - 8\tabcolsep) * \real{0.2250}}
  >{\centering\arraybackslash}p{(\linewidth - 8\tabcolsep) * \real{0.2625}}
  >{\raggedright\arraybackslash}p{(\linewidth - 8\tabcolsep) * \real{0.2500}}@{}}
\toprule\noalign{}
\begin{minipage}[b]{\linewidth}\centering
\(Q\)
\end{minipage} & \begin{minipage}[b]{\linewidth}\centering
\((c_1, c_2, c_3)\)
\end{minipage} & \begin{minipage}[b]{\linewidth}\centering
\(c_2 c_3\) (color)
\end{minipage} & \begin{minipage}[b]{\linewidth}\centering
\(c_1\) (weak isospin)
\end{minipage} & \begin{minipage}[b]{\linewidth}\raggedright
Correspondence
\end{minipage} \\
\midrule\noalign{}
\endhead
\bottomrule\noalign{}
\endlastfoot
0 & (0, 0, 0) & 00 (colorless) & 0 (down) & Lepton (neutral) \\
1 & (0, 0, 1) & 01 (color 1) & 0 (down) & Quark (down-type, color 1) \\
2 & (0, 1, 0) & 10 (color 2) & 0 (down) & Quark (down-type, color 2) \\
3 & (0, 1, 1) & 11 (color 3) & 0 (down) & Quark (down-type, color 3) \\
4 & (1, 0, 0) & 00 (colorless) & 1 (up) & Lepton (charged) \\
5 & (1, 0, 1) & 01 (color 1) & 1 (up) & Quark (up-type, color 1) \\
6 & (1, 1, 0) & 10 (color 2) & 1 (up) & Quark (up-type, color 2) \\
7 & (1, 1, 1) & 11 (color 3) & 1 (up) & Quark (up-type, color 3) \\
\end{longtable}
}

Particles with \(c_2 c_3 = 00\) correspond to leptons (no color charge),
and particles with \(c_2 c_3 \neq 00\) correspond to quarks (with color
charge). \(c_1\) encodes the up/down weak isospin.

Antiparticles correspond to the transformation \(Q \to (Q + 4) \bmod 8\)
(inversion of the \(c_1\) bit), accompanied by sign inversion of the
spatial axes.

\subsubsection{9.3 Correspondence to the 8 Gluon
States}\label{correspondence-to-the-8-gluon-states}

Of the 9 \(c_2 c_3\) transitions enumerated in §3, removing the linear
combination of one diagonal element (the state proportional to the color
singlet \(r\bar{r} + g\bar{g} + b\bar{b}\) in \(\mathrm{SU}(3)\)), the
\textbf{8 independent states are mapped to gluons \(g_1, \ldots, g_8\)}.

This reduction \(9 - 1 = 8\) originates from the adjoint representation
of the Lie algebra \(\mathfrak{su}(3)\) of \(\mathrm{SU}(3)\) (the
decomposition \(3 \otimes \bar{3} = 1 \oplus 8\)). While 9 transition
pairs are derived from the combinatorics of this model (§3), the
operation of removing 1 state to obtain 8 depends on the
representation-theoretic interpretation. The non-abelian structure
constants \(f^{abc}\) of \(\mathrm{SU}(3)\) are not derived from this
model.

\subsubsection{9.4 Correspondence to Gauge
Bosons}\label{correspondence-to-gauge-bosons}

Among the spin 1 sets enumerated in §4, the correspondence to Standard
Model gauge bosons is shown below.

{\def\LTcaptype{none} % do not increment counter
\begin{longtable}[]{@{}
  >{\raggedright\arraybackslash}p{(\linewidth - 4\tabcolsep) * \real{0.3000}}
  >{\raggedright\arraybackslash}p{(\linewidth - 4\tabcolsep) * \real{0.4333}}
  >{\centering\arraybackslash}p{(\linewidth - 4\tabcolsep) * \real{0.2667}}@{}}
\toprule\noalign{}
\begin{minipage}[b]{\linewidth}\raggedright
Fixed axis
\end{minipage} & \begin{minipage}[b]{\linewidth}\raggedright
Correspondence
\end{minipage} & \begin{minipage}[b]{\linewidth}\centering
Number of states
\end{minipage} \\
\midrule\noalign{}
\endhead
\bottomrule\noalign{}
\endlastfoot
\(t\) & \(W^+, W^-\) & 2 \\
\(R\) & \(\gamma, Z^0\) & 2 \\
\(Q\) (\(c_2 c_3\) transition, §9.3) & \(g_1, \ldots, g_8\) & 8 \\
\textbf{Total} & & \textbf{12} \\
\end{longtable}
}

The 2 states of the \(W\) boson (\(t = +1, -1\)) correspond to the
positive and negative electric charges, while \(\gamma\) and \(Z^0\)
correspond to the sign of the \(R\) axis (\(+\) vs \(-\)). The total of
12 matches the number of gauge boson states in the Standard Model:
\(\gamma + g \times 8 + W^+ + W^- + Z^0 = 12\).

The 8 states of \(c_1\) transitions (§3) can be interpreted as another
aspect of weak isospin transitions mediated by \(W\) bosons, but their
counting as independent particle states is treated in §9.10.

\subsubsection{9.5 Interpretation of Spin 2
Sets}\label{interpretation-of-spin-2-sets}

The following correspondences are assigned to the 3 types of spin 2 sets
enumerated in §5.

{\def\LTcaptype{none} % do not increment counter
\begin{longtable}[]{@{}
  >{\raggedright\arraybackslash}p{(\linewidth - 4\tabcolsep) * \real{0.1828}}
  >{\raggedright\arraybackslash}p{(\linewidth - 4\tabcolsep) * \real{0.1935}}
  >{\raggedright\arraybackslash}p{(\linewidth - 4\tabcolsep) * \real{0.6237}}@{}}
\toprule\noalign{}
\begin{minipage}[b]{\linewidth}\raggedright
Type
\end{minipage} & \begin{minipage}[b]{\linewidth}\raggedright
Fixed elements
\end{minipage} & \begin{minipage}[b]{\linewidth}\raggedright
Correspondence
\end{minipage} \\
\midrule\noalign{}
\endhead
\bottomrule\noalign{}
\endlastfoot
\(tR\) type & \(t, R\) & \textbf{Graviton} (correspondence within the
Standard Model + gravity framework) \\
\(t \cdot Q\)-transition type & \(t, Q\)-transition & No known particle
in the Standard Model \\
\(R \cdot Q\)-transition type & \(R, Q\)-transition & No known particle
in the Standard Model \\
\end{longtable}
}

The \(tR\) type is consistent with the properties of gravity (attraction
only). The \(tQ\) and \(RQ\) types are states whose existence is derived
from the combinatorial structure of this model, but they have no
counterpart in the current Standard Model.

\subsubsection{9.6 Interpretation of Spin 3
Sets}\label{interpretation-of-spin-3-sets}

The single type of spin 3 set (§6), the \(tRQ\)-transition type, has no
known corresponding particle in the Standard Model.

\subsubsection{9.7 Correspondence Between Sign Product and
Directionality of
Forces}\label{correspondence-between-sign-product-and-directionality-of-forces}

The following physical interpretation is assigned to the sign product
\(P\) defined in §7.

\begin{itemize}
\tightlist
\item
  \(P = +1\): The force is unidirectional (\textbf{attraction only})
\item
  \(P = -1\): The force is bidirectional (\textbf{both attraction and
  repulsion})
\end{itemize}

By Proposition 7.1, the \(tR\) type of spin 2 admits both \(P = +1\) and
\(P = -1\) configurations. Interpreting the \(P = +1\) configuration
(same sign for \(t, R\)) as mediating gravity yields attraction only,
which is consistent with the properties of gravity. What the \(P = -1\)
configuration (opposite signs for \(t, R\)) corresponds to physically is
unresolved, and remains as a future problem: either (a) a configuration
not physically realized, or (b) a correspondence to an unknown
phenomenon.

Similarly, the \(t\) and \(R\) types of spin 1 admit both \(P = \pm 1\)
configurations, consistent with forces having both attraction and
repulsion (electromagnetic force, weak force). The \(Q\)-transition type
of spin 1 has \(P = +1\) (unique), suggesting a relationship with the
confinement property of gluons.

\subsubsection{\texorpdfstring{9.8 Range of the \(R\) Axis and Coupling
Strength of
Forces}{9.8 Range of the R Axis and Coupling Strength of Forces}}\label{range-of-the-r-axis-and-coupling-strength-of-forces}

\textbf{Range of the \(R\) axis}: In prior work {[}W5{]}, the
number-theoretic conditions of central projection in discrete space lead
to the discretization \(R = \pm n^2 \cdot L_0\) (\(n \in \mathbb{N}\),
\(L_0\) is the unit length). The \(R\) axis connects the subjective
coordinate system with the spatial scale of the 6-dimensional structure,
and corresponds to the unboundedness of the mass spectrum.

\textbf{Interpretation of the correspondence between each force and the
axes of this model}:

{\def\LTcaptype{none} % do not increment counter
\begin{longtable}[]{@{}
  >{\raggedright\arraybackslash}p{(\linewidth - 6\tabcolsep) * \real{0.1402}}
  >{\raggedright\arraybackslash}p{(\linewidth - 6\tabcolsep) * \real{0.1215}}
  >{\raggedright\arraybackslash}p{(\linewidth - 6\tabcolsep) * \real{0.5140}}
  >{\raggedright\arraybackslash}p{(\linewidth - 6\tabcolsep) * \real{0.2243}}@{}}
\toprule\noalign{}
\begin{minipage}[b]{\linewidth}\raggedright
Force
\end{minipage} & \begin{minipage}[b]{\linewidth}\raggedright
Mediator
\end{minipage} & \begin{minipage}[b]{\linewidth}\raggedright
Involved axes
\end{minipage} & \begin{minipage}[b]{\linewidth}\raggedright
Coupling characteristics
\end{minipage} \\
\midrule\noalign{}
\endhead
\bottomrule\noalign{}
\endlastfoot
Strong force & \(g_1\)--\(g_8\) & \(Q\) axis (\(c_2 c_3\) color
transitions, 8 channels, §9.3) & Only between particles with color
charge \\
Weak force & \(W^\pm, Z^0\) & \(t\) axis + \(Q\) axis (\(c_1\)
transitions) & All fermions \\
Electromagnetic & \(\gamma\) & \(R\) axis & Between particles with
nonzero electric charge \\
Gravity & \(G\) & \(t, R\), both axes nonzero (\(n = 2\)) & All
particles (\(Q\)-independent) \\
\end{longtable}
}

\subsubsection{9.9 Table A: Correspondence Between Confirmed Particles
and 6-Dimensional
Sets}\label{table-a-correspondence-between-confirmed-particles-and-6-dimensional-sets}

This table shows particles only. Antifermions are obtained by sign
inversion of spatial axes and \(Q \to (Q + 4) \bmod 8\), yielding 24
states: 6 antileptons + 18 antiquarks. For bosons, particle-antiparticle
pairs such as \(W^+/W^-\) are explicitly listed in the table;
\(\gamma, Z^0, H^0, G\), and gluons are self-conjugate.

\textbf{Bosons}:

{\def\LTcaptype{none} % do not increment counter
\begin{longtable}[]{@{}
  >{\raggedright\arraybackslash}p{(\linewidth - 10\tabcolsep) * \real{0.1408}}
  >{\raggedright\arraybackslash}p{(\linewidth - 10\tabcolsep) * \real{0.1549}}
  >{\centering\arraybackslash}p{(\linewidth - 10\tabcolsep) * \real{0.0704}}
  >{\raggedright\arraybackslash}p{(\linewidth - 10\tabcolsep) * \real{0.4225}}
  >{\centering\arraybackslash}p{(\linewidth - 10\tabcolsep) * \real{0.0704}}
  >{\centering\arraybackslash}p{(\linewidth - 10\tabcolsep) * \real{0.1408}}@{}}
\toprule\noalign{}
\begin{minipage}[b]{\linewidth}\raggedright
Particle
\end{minipage} & \begin{minipage}[b]{\linewidth}\raggedright
Symbol
\end{minipage} & \begin{minipage}[b]{\linewidth}\centering
\(s\)
\end{minipage} & \begin{minipage}[b]{\linewidth}\raggedright
Set \((x, y, z, t, R, Q)\)
\end{minipage} & \begin{minipage}[b]{\linewidth}\centering
\(n\)
\end{minipage} & \begin{minipage}[b]{\linewidth}\centering
Mass (MeV)
\end{minipage} \\
\midrule\noalign{}
\endhead
\bottomrule\noalign{}
\endlastfoot
Higgs & \(H^0\) & 0 & \((0, 0, +, 0, 0, 0)\) & 0 & 125,250 \\
\(W\) boson & \(W^+\) & 1 & \((0, 0, 0, +, 0, 0)\) & 1 & 80,379 \\
\(W\) boson & \(W^-\) & 1 & \((0, 0, 0, -, 0, 0)\) & 1 & 80,379 \\
Photon & \(\gamma\) & 1 & \((0, 0, 0, 0, +, 0)\) & 1 & 0 \\
\(Z\) boson & \(Z^0\) & 1 & \((0, 0, 0, 0, -, 0)\) & 1 & 91,188 \\
Gluon & \(g_1\)--\(g_8\) & 1 & \((0, 0, 0, 0, 0, Q_i \to Q_j)\) & 1 &
0 \\
Graviton & \(G\) & 2 & \((0, 0, 0, \pm, \pm, 0)\) & 2 & 0 \\
\end{longtable}
}

\textbf{Verification}: Higgs: \(z\) axis only nonzero, \(k = 1\)
(bosonic type), \(n = 0\), \(s = 0\) \(\checkmark\). \(W\): \(k = 1\)
(bosonic type), \(t\) nonzero so \(n = 1\), \(s = 1\) \(\checkmark\).
Photon: \(k = 0\) (bosonic type), \(R\) nonzero so \(n = 1\), \(s = 1\)
\(\checkmark\). Gluon: \(k = 0\) (bosonic type), \(Q\)-transition so
\(n = 1\), \(s = 1\) \(\checkmark\). Graviton: \(k = 0\) (bosonic type),
\(t, R\) nonzero so \(n = 2\), \(s = 2\) \(\checkmark\).

\textbf{Selection of the \(z\) axis for the Higgs}: The nonzero axis of
the Higgs boson involves an interpretive degree of freedom in choosing 1
axis from the 4 axes \(xyzt\). For consistency with the generation
labels (generation 1: \(x\), generation 2: \(y\), generation 3: \(z\)),
the \(z\) axis (3rd generation) is selected. This choice is consistent
with the fact that 3rd-generation fermions (top: \(172{,}760\) MeV) have
the largest mass, provided the coupling strength depends on the overlap
along shared spatial axes (see §9.11.7).

\textbf{Leptons}:

{\def\LTcaptype{none} % do not increment counter
\begin{longtable}[]{@{}
  >{\raggedright\arraybackslash}p{(\linewidth - 12\tabcolsep) * \real{0.2346}}
  >{\raggedright\arraybackslash}p{(\linewidth - 12\tabcolsep) * \real{0.1235}}
  >{\centering\arraybackslash}p{(\linewidth - 12\tabcolsep) * \real{0.0617}}
  >{\centering\arraybackslash}p{(\linewidth - 12\tabcolsep) * \real{0.1235}}
  >{\raggedright\arraybackslash}p{(\linewidth - 12\tabcolsep) * \real{0.2593}}
  >{\centering\arraybackslash}p{(\linewidth - 12\tabcolsep) * \real{0.0617}}
  >{\raggedright\arraybackslash}p{(\linewidth - 12\tabcolsep) * \real{0.1358}}@{}}
\toprule\noalign{}
\begin{minipage}[b]{\linewidth}\raggedright
Particle
\end{minipage} & \begin{minipage}[b]{\linewidth}\raggedright
Symbol
\end{minipage} & \begin{minipage}[b]{\linewidth}\centering
\(s\)
\end{minipage} & \begin{minipage}[b]{\linewidth}\centering
Generation
\end{minipage} & \begin{minipage}[b]{\linewidth}\raggedright
Set
\end{minipage} & \begin{minipage}[b]{\linewidth}\centering
\(Q\)
\end{minipage} & \begin{minipage}[b]{\linewidth}\raggedright
Mass (MeV)
\end{minipage} \\
\midrule\noalign{}
\endhead
\bottomrule\noalign{}
\endlastfoot
Electron & \(e^-\) & 1/2 & 1 & \((+, 0, 0, +, 0, 4)\) & 4 & 0.511 \\
Electron neutrino & \(\nu_e\) & 1/2 & 1 & \((+, +, 0, 0, 0, 0)\) & 0 &
\(< 10^{-6}\) \\
Muon & \(\mu^-\) & 1/2 & 2 & \((0, +, 0, +, 0, 4)\) & 4 & 105.66 \\
Muon neutrino & \(\nu_\mu\) & 1/2 & 2 & \((0, +, +, 0, 0, 0)\) & 0 &
\(< 0.17\) \\
Tau & \(\tau^-\) & 1/2 & 3 & \((0, 0, +, +, 0, 4)\) & 4 & 1,776.9 \\
Tau neutrino & \(\nu_\tau\) & 1/2 & 3 & \((+, 0, +, 0, 0, 0)\) & 0 &
\(< 15.5\) \\
\end{longtable}
}

\textbf{Verification}: Electron: \(x, t\) are nonzero among \(xyzt\), so
\(k = 2\) (fermionic type); \(\{R, Q\text{-transition}\}\) are both not
applicable (\(R = 0\), \(Q = 4\) is a static label, not a transition),
so \(n = 0\), \(s = 0 + 1/2 = 1/2\) \(\checkmark\). Neutrino: \(x, y\)
are nonzero, so \(k = 2\) (fermionic type), \(n = 0\), \(s = 1/2\)
\(\checkmark\).

\textbf{Characteristics of charged leptons}: The \(t\) axis is nonzero,
with \(Q = 4\) (\(c_1 = 1, c_2 c_3 = 00\); weak isospin up, colorless).
Neutrinos have \(t = 0\), \(Q = 0\), with only 2 spatial axes among
\(xyzt\) being nonzero.

\textbf{Distinction of generations}: For charged leptons, 1 nonzero
spatial axis and the \(t\) axis satisfy \(k = 2\); for neutrinos, 2
spatial axes satisfy \(k = 2\). The generation label is determined by
the nonzero spatial axes (generation 1: \(x\), generation 2: \(y\),
generation 3: \(z\)).

\textbf{Generation correspondence rule}: In the lepton doublet
\(\{\ell^-_n, \nu_n\}\) of the same generation, one of the 2 nonzero
spatial axes of \(\nu_n\) coincides with the nonzero spatial axis of
\(\ell^-_n\). The choice of the remaining axis is an interpretive degree
of freedom; this paper adopts a cyclic assignment (generation 1: \(xy\),
generation 2: \(yz\), generation 3: \(xz\)).

Antileptons correspond to sign inversion of spatial axes and
\(Q \to (Q + 4) \bmod 8\) (\(e^+\): \(Q = 4 \to 0\)).

\textbf{Quarks}:

{\def\LTcaptype{none} % do not increment counter
\begin{longtable}[]{@{}
  >{\raggedright\arraybackslash}p{(\linewidth - 12\tabcolsep) * \real{0.1127}}
  >{\raggedright\arraybackslash}p{(\linewidth - 12\tabcolsep) * \real{0.0845}}
  >{\centering\arraybackslash}p{(\linewidth - 12\tabcolsep) * \real{0.0704}}
  >{\centering\arraybackslash}p{(\linewidth - 12\tabcolsep) * \real{0.1408}}
  >{\raggedright\arraybackslash}p{(\linewidth - 12\tabcolsep) * \real{0.3803}}
  >{\centering\arraybackslash}p{(\linewidth - 12\tabcolsep) * \real{0.0704}}
  >{\raggedright\arraybackslash}p{(\linewidth - 12\tabcolsep) * \real{0.1408}}@{}}
\toprule\noalign{}
\begin{minipage}[b]{\linewidth}\raggedright
Particle
\end{minipage} & \begin{minipage}[b]{\linewidth}\raggedright
Symbol
\end{minipage} & \begin{minipage}[b]{\linewidth}\centering
\(s\)
\end{minipage} & \begin{minipage}[b]{\linewidth}\centering
Generation
\end{minipage} & \begin{minipage}[b]{\linewidth}\raggedright
Set
\end{minipage} & \begin{minipage}[b]{\linewidth}\centering
\(Q\)
\end{minipage} & \begin{minipage}[b]{\linewidth}\raggedright
Mass (MeV)
\end{minipage} \\
\midrule\noalign{}
\endhead
\bottomrule\noalign{}
\endlastfoot
Up & \(u\) & 1/2 & 1 & \((+, 0, 0, +, 0, 5/6/7)\) & 5,6,7 & 2.16 \\
Down & \(d\) & 1/2 & 1 & \((+, 0, 0, -, 0, 1/2/3)\) & 1,2,3 & 4.67 \\
Charm & \(c\) & 1/2 & 2 & \((0, +, 0, +, 0, 5/6/7)\) & 5,6,7 & 1,270 \\
Strange & \(s\) & 1/2 & 2 & \((0, +, 0, -, 0, 1/2/3)\) & 1,2,3 & 93.4 \\
Top & \(t\) & 1/2 & 3 & \((0, 0, +, +, 0, 5/6/7)\) & 5,6,7 & 172,760 \\
Bottom & \(b\) & 1/2 & 3 & \((0, 0, +, -, 0, 1/2/3)\) & 1,2,3 & 4,180 \\
\end{longtable}
}

\textbf{Verification}: Up quark: \(x, t\) are nonzero among \(xyzt\), so
\(k = 2\) (fermionic type); \(Q = 5\) is a static label (not a
transition), so \(n = 0\), \(s = 0 + 1/2 = 1/2\) \(\checkmark\).

\textbf{Characteristics of quarks}: All quarks have the \(t\) axis
nonzero. Up-type quarks (\(u, c, t\)) have \(t > 0\) with
\(Q \in \{5, 6, 7\}\) (\(c_1 = 1\); weak isospin up, colored). Down-type
quarks (\(d, s, b\)) have \(t < 0\) with \(Q \in \{1, 2, 3\}\)
(\(c_1 = 0\); weak isospin down, colored). The sign of the \(t\) axis
reflects the sign structure of the electric charge.

Each quark has 3 color states corresponding to the color values of \(Q\)
(\(c_2 c_3 \in \{01, 10, 11\}\)). Antiquarks correspond to sign
inversion of spatial axes and \(Q \to (Q + 4) \bmod 8\).

\textbf{Summary}:

{\def\LTcaptype{none} % do not increment counter
\begin{longtable}[]{@{}
  >{\raggedright\arraybackslash}p{(\linewidth - 8\tabcolsep) * \real{0.2712}}
  >{\centering\arraybackslash}p{(\linewidth - 8\tabcolsep) * \real{0.0424}}
  >{\centering\arraybackslash}p{(\linewidth - 8\tabcolsep) * \real{0.1356}}
  >{\centering\arraybackslash}p{(\linewidth - 8\tabcolsep) * \real{0.2373}}
  >{\centering\arraybackslash}p{(\linewidth - 8\tabcolsep) * \real{0.3136}}@{}}
\toprule\noalign{}
\begin{minipage}[b]{\linewidth}\raggedright
Classification
\end{minipage} & \begin{minipage}[b]{\linewidth}\centering
\(s\)
\end{minipage} & \begin{minipage}[b]{\linewidth}\centering
\(Q\)
\end{minipage} & \begin{minipage}[b]{\linewidth}\centering
States derived in this model
\end{minipage} & \begin{minipage}[b]{\linewidth}\centering
SM correspondence
\end{minipage} \\
\midrule\noalign{}
\endhead
\bottomrule\noalign{}
\endlastfoot
Higgs & 0 & 0 & 1 & 1 \\
Gauge boson (\(t\) type) & 1 & 0 & 2 & \(W^+, W^-\) \\
Gauge boson (\(R\) type) & 1 & 0 & 2 & \(\gamma, Z^0\) \\
Gluon (\(Q\)-transition type) & 1 & Transition & 8 &
\(g_1, \ldots, g_8\) \\
Graviton (\(tR\) type) & 2 & 0 & 1 & \(G\) (beyond SM) \\
Lepton & 1/2 & 0 or 4 & 6 & 6 \\
Antilepton & 1/2 & 0 or 4 & 6 & 6 \\
Quark & 1/2 & 1,2,3 or 5,6,7 & 18 & 18 \\
Antiquark & 1/2 & 1,2,3 or 5,6,7 & 18 & 18 \\
\textbf{Subtotal} & & & \textbf{62} & \textbf{62 (61 + G)} \\
\end{longtable}
}

The 62 states map to 61 Standard Model states + 1 graviton. The excess
color singlet state in the previous model (\(Q = 4\) values) has been
resolved by the extension to \(Q = 8\) values.

In addition to the above, combinations with \(s = 3/2, 2, 5/2, 3\)
(Table B, §9.10) are derived from this model, but no known particles in
the Standard Model correspond to them.

\subsubsection{9.10 Table B: Combinations with No Counterpart in the
Standard
Model}\label{table-b-combinations-with-no-counterpart-in-the-standard-model}

Listed below are the states derived from the combinatorial structure of
this model for which no known particle in the Standard Model exists.
These are combinatorial consequences within this model and do not assert
the existence of undiscovered particles.

\textbf{Bosons (number of nonzero axes among \(xyzt\): \(k \neq 2\))}:

{\def\LTcaptype{none} % do not increment counter
\begin{longtable}[]{@{}
  >{\raggedright\arraybackslash}p{(\linewidth - 8\tabcolsep) * \real{0.1023}}
  >{\centering\arraybackslash}p{(\linewidth - 8\tabcolsep) * \real{0.0568}}
  >{\centering\arraybackslash}p{(\linewidth - 8\tabcolsep) * \real{0.0568}}
  >{\raggedright\arraybackslash}p{(\linewidth - 8\tabcolsep) * \real{0.4545}}
  >{\centering\arraybackslash}p{(\linewidth - 8\tabcolsep) * \real{0.3295}}@{}}
\toprule\noalign{}
\begin{minipage}[b]{\linewidth}\raggedright
Type
\end{minipage} & \begin{minipage}[b]{\linewidth}\centering
\(s\)
\end{minipage} & \begin{minipage}[b]{\linewidth}\centering
\(n\)
\end{minipage} & \begin{minipage}[b]{\linewidth}\raggedright
Set form
\end{minipage} & \begin{minipage}[b]{\linewidth}\centering
Number of states
\end{minipage} \\
\midrule\noalign{}
\endhead
\bottomrule\noalign{}
\endlastfoot
\(tQ\) type & 2 & 2 & \((0, 0, 0, \pm, 0, Q_i \to Q_j)\) &
\(2 \times N_Q\) \\
\(RQ\) type & 2 & 2 & \((0, 0, 0, 0, \pm, Q_i \to Q_j)\) &
\(2 \times N_Q\) \\
\(tRQ\) type & 3 & 3 & \((0, 0, 0, \pm, \pm, Q_i \to Q_j)\) &
\(4 \times N_Q\) \\
\end{longtable}
}

Here \(N_Q\) is the number of \(Q\) transition states: \(c_2 c_3\)
transition 9 states + \(c_1\) transition 8 states = 17 (§3). Under the
interpretation of §9.3 where \(c_2 c_3\) transitions are reduced
\(9 \to 8\), \(N_Q = 16\).

\textbf{Fermions (number of nonzero axes among \(xyzt\): \(k = 2\),
excluding combinations included in Table A)}:

With the extension to \(Q = 8\) values, the number of fermionic-type
combinations increases substantially. Choosing 2 axes from \(xyzt\):
\(\binom{4}{2} = 6\) ways \(\times\) signs of each axis \(2^2 = 4\) ways
\(\times\) 8 \(Q\) values \(= 6 \times 4 \times 8 = 192\) fermionic
states exist. Of these, the 48 states mapped in Table A (6 leptons + 6
antileptons + 18 quarks + 18 antiquarks) are subtracted, leaving the
remainder as unassigned states.

Higher-spin fermions (with \(R\) axis or \(Q\)-transition nonzero,
\(n \geq 1\)):

{\def\LTcaptype{none} % do not increment counter
\begin{longtable}[]{@{}lccl@{}}
\toprule\noalign{}
Type & \(s\) & \(n\) & Additional nonzero axis \\
\midrule\noalign{}
\endhead
\bottomrule\noalign{}
\endlastfoot
\(R\) type & 3/2 & 1 & \(R\) \\
\(Q\)-transition type & 3/2 & 1 & \(Q\) (transition) \\
\(RQ\) type & 5/2 & 2 & \(R, Q\) (transition) \\
\end{longtable}
}

These higher-spin fermions are states where, in addition to 2 nonzero
axes among \(xyzt\), \(R\) or \(Q\)-transitions are also nonzero. They
have an origin distinct from the partner particles of supersymmetric
theories.

\subsubsection{9.11 New Structural Consequences and Open Problems in the
Interpretation
Example}\label{new-structural-consequences-and-open-problems-in-the-interpretation-example}

\textbf{9.11.1 Initial symmetry and value structure of the \(xyzt\)
axes}

Since the 4 axes \(xyzt\) are completely equivalent in their initial
symmetry, each axis is expected to take values of the same kind
(positive integers). However, depending on the number of nonzero axes,
there exist 4 types of values \(v_1, v_2, v_3, v_4\) (for the cases of
1, 2, 3, and 4 nonzero axes, respectively), and these are generally
distinct.

What is structurally established is only that
\(v_1 \neq v_2 \neq v_3 \neq v_4\) is permitted. Everything below is
\textbf{interpretive speculation} and is not derived from within this
model.

\textbf{Relation to dimensional analysis}: In the central projection
model of prior work {[}W5{]}, spatial length \(L\) and mass \(M\) have
the dimensional relation \(M = L\) (in natural units). Applying this
relation to this model, the sum of absolute values of coordinate values
along the \(xyzt\) axes may be related to the mass metric \(M(\sigma)\).
That is, for a set \(\sigma = (x, y, z, t, R, Q)\),

\[M(\sigma) \sim \sum_{i \in \{x, y, z, t\}} |x_i|\]

may serve as the origin of mass. However, how the \(R\) axis (spatial
scale axis) and \(Q\) axis (gauge symmetry axis) contribute to this
metric is undetermined, and a more general form including \(R, Q\),

\[M(\sigma) \sim f\!\left(\sum |x_i|,\, R,\, Q\right)\]

cannot be excluded.

\textbf{Speculation that \(v_1 \gg v_2\)}: While the mass of the Higgs
boson (\(k = 1\), \(z\) axis only nonzero, \(v_1\)) is \(125{,}250\)
MeV, the majority of fermions (\(k = 2\), 2 nonzero axes, \(v_2\)) are
far lighter (electron \(0.511\) MeV, neutrinos \(< 10^{-6}\) MeV). If
the above dimensional analysis speculation is correct, then
\(v_1 \gg v_2\) must hold, suggesting that the coordinate value for the
case of only 1 nonzero axis is substantially larger than that for 2
nonzero axes.

\textbf{Inter-generation mass hierarchy}: The assignment of the Higgs
boson's nonzero axis to \(z\) (the 3rd-generation axis) provides a
structural indication for the mass differences among generations within
the \(v_2\) range. 3rd-generation fermions (\(z\) axis nonzero) directly
share a spatial axis with the Higgs standing wave and therefore have the
strongest coupling, while 1st-generation (\(x\) axis) and 2nd-generation
(\(y\) axis) fermions lack a shared axis and are limited to indirect
coupling. This is consistent with the empirical hierarchy
\(m_3 \gg m_{1,2}\) (top \(172{,}760\) MeV \(\gg\) charm \(1{,}270\)
MeV, up \(2.16\) MeV). However, since the \(x\) and \(y\) axes are
locally equivalent under \(\mathrm{SO}(3)\) symmetry, the difference
\(m_2 \neq m_1\) cannot be derived within this framework. The
quantitative derivation of mass ratios remains an open problem.

\textbf{Relation to the problem of undiscovered \(k = 3, 4\) states}:
States with 3 or 4 nonzero axes among \(xyzt\) (\(k = 3, 4\)) are
derived from this model as bosonic type (\(k \neq 2\)), but no known
particles in Table A correspond to them. If the values \(v_3, v_4\)
corresponding to these states are extremely small (or near zero), the
energy scale of the corresponding particles may exceed the current
experimental reach. Conversely, if \(v_3, v_4\) are extremely large, the
corresponding particles may not be physically realized due to stability
issues. In either case, the fact that \(k = 3, 4\) states remain
undiscovered is likely closely related to the value structure of the
\(xyzt\) axes.

\textbf{9.11.2 The \(Q = 0\) constraint and free particles}

\(Q = 0\) signifies neutrality on the gauge symmetry axis. As an
interpretation, the constraint that only particles with \(Q = 0\) are
observable as free particles is conceivable. This can be interpreted as
the geometric origin of quark confinement (particles with \(Q \neq 0\)
are not observed individually).

\textbf{9.11.3 Predictions of undiscovered particles}

With the extension of the \(Q\) axis from 4 to 8 values, the number of
combinatorial states derived from this model has increased
substantially. These combinations beyond the 62 states of Table A are to
be examined in future research as either (a) predictions of undiscovered
particles, (b) physically unrealized combinations (existence of
selection rules), or (c) limitations of this interpretation.

In particular, higher-spin fermions (\(s = 3/2, 5/2\)) and additional
bosons (\(tQ\) type, \(RQ\) type spin 2; \(tRQ\) type spin 3) have no
counterparts among known particles.

\textbf{9.11.4 Derivation of mass ratios}

Deriving the known particle mass ratios from the set-theoretic structure
of this model is a task for future work. The specific ratios of the 4
types of values \(v_1, v_2, v_3, v_4\) for the \(xyzt\) axes (§9.11.1),
and the form of the contribution of the \(R\) and \(Q\) axes to the mass
metric \(M(\sigma)\), may be key to this problem. In particular, if the
relation \(v_1 \gg v_2\) (§9.11.1) is confirmed, a geometric explanation
for the mass hierarchy between the Higgs boson and fermions would be
obtained.

\textbf{9.11.5 Derivation of gauge groups}

A correspondence is suggested between the 8 values (3-bit structure) of
the \(Q\) axis and \(\mathrm{SU}(3)_C \times \mathrm{SU}(2)_L\), and
between the continuous scale of the \(R\) axis and \(\mathrm{U}(1)_Y\),
but the rigorous mechanism for the transition from discrete structure to
continuous gauge groups remains unproven.

\textbf{9.11.6 CKM matrix and PMNS matrix}

How generation mixing (CKM matrix, PMNS matrix) is derived as the
breaking of \(S_3\) symmetry of the \(xyz\) axes is an open problem.

\textbf{9.11.7 Higgs mechanism and generation mass hierarchy}

The Higgs boson corresponds to only the \(z\) axis being nonzero
(\(k = 1\), Table A: \((0, 0, +, 0, 0, 0)\)). The \(z\) axis is the
label axis for the 3rd generation (top, bottom, tau), and this
assignment has the following consistency:

\begin{itemize}
\tightlist
\item
  Since the coupling strength to the Higgs depends on the overlap along
  shared spatial axes, 3rd-generation fermions with nonzero \(z\) axis
  have the strongest coupling and acquire the largest mass (top:
  \(172{,}760\) MeV).
\item
  1st-generation (\(x\) axis) and 2nd-generation (\(y\) axis) fermions
  do not share the Higgs's primary oscillation axis, resulting in weaker
  coupling and lighter mass.
\end{itemize}

How spontaneous symmetry breaking by the Higgs field is understood
within this framework remains a task for future work, but the selection
of the \(z\) axis provides a geometric explanation for the
inter-generation mass hierarchy (\(m_3 \gg m_{1,2}\)).

\textbf{9.11.8 Complete derivation of charge values}

The sign of the electric charge can be read from the sign of the \(t\)
axis, and the presence/absence of color charge and the absolute value of
electric charge (1 vs 1/3) from \(c_2 c_3\) of the \(Q\) axis, but the
complete derivation of the asymmetry between \(+2/3\) and \(-1/3\) is a
task for future work.

\subsubsection{9.12 Conclusion of the Interpretation
Example}\label{conclusion-of-the-interpretation-example}

This section presented one example of mapping the mathematical
properties of the 6-dimensional structure defined and derived in §1--§8
to the particle classification of the Standard Model.

\begin{enumerate}
\def\labelenumi{\arabic{enumi}.}
\tightlist
\item
  The 8 values (3-bit structure) of the \(Q\) axis discretely encode the
  \(\mathrm{SU}(3)_C \times \mathrm{SU}(2)_L\) gauge symmetry.
\item
  From the initial symmetry of the 4 axes \(xyzt\), the definition of
  spin arises naturally (fermion: \(k = 2\); boson: \(k \neq 2\)).
\item
  62 states can be mapped to known particles of the Standard Model (61 +
  graviton).
\item
  Of the 9 \(c_2 c_3\) transitions, the 8 states remaining after
  removing the color singlet correspond to gluons (the reduction
  \(9 \to 8\) depends on \(\mathrm{SU}(3)\) representation theory).
\item
  The 12 states of spin 1 sets correspond to the gauge bosons of the
  Standard Model.
\item
  Charged fermions have the \(t\) axis nonzero, and the sign of the
  electric charge corresponds to the sign of \(t\).
\item
  The \(Q = 0\) constraint can be interpreted as the geometric origin of
  quark confinement.
\item
  The value structure of the \(xyzt\) axes (4 types of positive integer
  values) may contain the origin of the mass hierarchy.
\item
  Higher-spin states and surplus states accompanying the extension to
  \(Q = 8\) values suggest either predictions of undiscovered particles
  or the existence of selection rules.
\end{enumerate}

The derivation of specific mass values, charge values, the non-abelian
structure of gauge groups, and generation mixing matrices are tasks for
future work (§9.11).

\begin{center}\rule{0.5\linewidth}{0.5pt}\end{center}

\subsection{References}\label{references}

{[}1{]} Finkelstein, D. \& Rubinstein, J. ``Connection between Spin,
Statistics, and Kinks.'' \emph{J. Math. Phys.} 9, 1762 (1968).

{[}2{]} Atiyah, M.F., Manton, N.S. \& Schroers, B.J. ``Geometric Models
of Matter.'' \emph{Proc. R. Soc. A} 468, 1252 (2012). arXiv: 1108.5151.

{[}3{]} Furey, C. ``Standard Model Physics from an Algebra?'' arXiv:
1611.09182 (2016).

{[}4{]} Kaluza, Th. ``Zum Unitätsproblem der Physik.''
\emph{Sitzungsber. Preuss. Akad. Wiss. Berlin} 966--972 (1921).

{[}5{]} Coxeter, H.S.M. \emph{Regular Polytopes}. Dover (1973).

{[}6{]} Kihara, N. ``Bivectorial Classification of Spin Derived from the
Orientation Structure of the 5-Dimensional Hypercube.'' DOI:
\href{https://doi.org/10.5281/zenodo.19643358}{10.5281/zenodo.19643358}
(2026).

{[}7{]} Kihara, N. ``Geometric Classification of Spin Derived from the
Orientation Structure of the 5-Dimensional Orthoplex.'' DOI:
\href{https://doi.org/10.5281/zenodo.19630972}{10.5281/zenodo.19630972}
(2026).

{[}W5{]} Kihara, N. ``Complete Spherical Coverage of Central Projection
in Discrete Space and the Number-Theoretic Necessity of the
5-Dimensional Background Space.'' DOI:
\href{https://doi.org/10.5281/zenodo.19624957}{10.5281/zenodo.19624957}
(2026).
\end{document}
